% ============================================================
% 第9章 有机化合物的结构特点与烃（必修二第七章、选必3第一、二章）
% ============================================================
\chapter{有机物的结构与烃}

\begin{daoxue}
\textbf{知识主线：}有机物为什么种类繁多（碳的成键特点、同分异构）$\to$ 怎样分类和命名（官能团、系统命名）$\to$ 怎样确定结构（分离提纯、元素分析、质谱、红外、核磁）$\to$ 最简单的有机物——烃的性质（烷：取代；烯、炔：加成、氧化、加聚；芳香烃：取代为主）。

\textbf{高考考情：}广东卷第20题（有机大题，14分）每年必考官能团名称、分子式、同分异构体数目或结构、核磁共振氢谱、碳原子杂化与共面（2024年第20题(2)、2025年第20题(2)(4)、2026年第20题(2)(3)）；选择题第9题左右考陌生有机物的结构与性质（2024年第9题柳珊瑚酸）。

\textbf{学习目标：}①会判断官能团、碳原子杂化方式和共线共面；②会有序书写同分异构体并计数；③会用谱图信息推断结构；④掌握各类烃的典型反应。
\end{daoxue}

\section{有机物的结构特点}

\subsection{碳原子的成键特点与官能团}
碳原子最外层 4 个电子，一般形成 4 个共价键；碳原子之间可形成单键、双键、三键，可以成链也可以成环——这是有机物种类繁多的原因。决定有机物特殊性质的原子或原子团叫\textbf{官能团}：

\begin{center}
\small
\begin{tabular}{llll}
\toprule
类别 & 官能团 & 类别 & 官能团 \\
\midrule
烯烃 & 碳碳双键 \ce{C=C} & 醛 & 醛基 \ce{-CHO} \\
炔烃 & 碳碳三键 \ce{C#C} & 酮 & 酮羰基 \ce{C=O} \\
卤代烃 & 碳卤键 \ce{-X} & 羧酸 & 羧基 \ce{-COOH} \\
醇 & （醇）羟基 \ce{-OH}（连在饱和碳上） & 酯 & 酯基 \ce{-COOR} \\
酚 & （酚）羟基 \ce{-OH}（直接连苯环） & 胺 & 氨基 \ce{-NH2} \\
醚 & 醚键 \ce{-O-} & 酰胺 & 酰胺基 \ce{-CONH-} \\
\bottomrule
\end{tabular}
\end{center}

\subsection{碳原子的杂化与分子的空间结构}
\textbf{饱和碳}（4 个单键）为 $\mathrm{sp}^3$ 杂化，与相连原子呈四面体；\textbf{双键碳}（含苯环碳、羰基碳）为 $\mathrm{sp}^2$ 杂化，与相连原子共平面；\textbf{三键碳}为 sp 杂化，与相连原子共直线。

\begin{fangfa}[共线共面的判断]
以三种基本模型为“积木”：\textbf{甲烷}（四面体，最多 3 个原子共面）；\textbf{乙烯}（6 个原子共面）；\textbf{乙炔}（4 个原子共线）；\textbf{苯}（12 个原子共面，对位上的 4 个原子共线）。单键可以旋转，所以“可能共面”与“一定共面”要区分：只要分子中有饱和碳原子（如 \ce{-CH3}、\ce{-CH2-}），所有原子就\textbf{不可能}共面。2024年广东卷第9题中，柳珊瑚酸的环上有一个碳同时与 4 个碳相连，为 $\mathrm{sp}^3$ 杂化（四面体），因此“环系结构中 3 个五元环共平面”是错误的。
\end{fangfa}

\subsection{同分异构现象}
分子式相同而结构不同的化合物互为同分异构体。类型：①\textbf{碳架异构}（如正丁烷与异丁烷）；②\textbf{位置异构}（如 1-丁烯与 2-丁烯）；③\textbf{官能团异构}（如乙醇与二甲醚，丙醛与丙酮，羧酸与酯）；④\textbf{顺反异构}（双键两端碳原子各连两个不同基团时，如 2-丁烯有顺、反两种）；⑤\textbf{对映异构}（含手性碳原子）。

\begin{fangfa}[同分异构体的有序书写与计数]
\textbf{①等效氢法求一元取代物数目：}同一碳上的 H 等效；同一碳上所连甲基上的 H 等效；处于对称位置的 H 等效。一元取代物种数 $=$ 等效氢种数。如 \ce{(CH3)3CCH2CH3} 有 3 种等效氢。\\
\textbf{②基元法：}丙基有 2 种、丁基有 4 种、戊基有 8 种。所以 \ce{C4H9Cl} 有 4 种，\ce{C4H9OH}（醇）有 4 种，\ce{C4H9CHO}（醛）、\ce{C4H9COOH}（羧酸）各有 4 种，\ce{C5H11OH}（醇）有 8 种。\\
\textbf{③定一移一法求二元取代物：}先固定一个取代基，再移动另一个。\\
\textbf{④苯环上的取代：}两个取代基有邻、间、对 3 种；三个取代基\textbf{都相同}有 3 种，\textbf{两同一异}有 6 种，\textbf{三个都不同}有 10 种。\\
\textbf{⑤酯类：}按“酸的碳数$+$醇的碳数”拆分，如 \ce{C4H8O2} 的酯：甲酸丙酯（2 种）、乙酸乙酯（1 种）、丙酸甲酯（1 种），共 4 种。
\end{fangfa}

\begin{zhenti}{2025年广东卷第20题(2)②\quad 来源：仓库所附真题文档}
化合物 2a（苯乙烯，\ce{C8H8}）可与 \ce{H2O} 发生加成反应生成化合物Ⅰ。在Ⅰ的同分异构体中，同时含有苯环和醇羟基结构的共\tk{}种（含化合物Ⅰ，不考虑立体异构）。

\textbf{解析：}Ⅰ的分子式为 \ce{C8H10O}，除苯环外剩 2 个 C，均为饱和结构。醇羟基不能直接连在苯环上：①一个侧链：\ce{C6H5CH2CH2OH}、\ce{C6H5CH(OH)CH3}，2 种；②两个侧链 \ce{-CH3} 和 \ce{-CH2OH}：邻、间、对 3 种。共 5 种。\\
\dianbo “醇羟基”排除了酚（\ce{-OH} 直接连苯环）；同时要排除醚（如 \ce{C6H5CH2OCH3}）。
\end{zhenti}

\begin{liti}[限定条件的同分异构体]
分子式为 \ce{C8H8O2}、含苯环且属于酯类的同分异构体共有几种？其中核磁共振氢谱有 4 组峰的是哪些？
\tcblower
\silu 不饱和度 $\Omega=\frac{2\times8+2-8}{2}=5$：苯环 4，酯基 1，其余均饱和。按“酸 $+$ 醇（酚）”拆分。\\
\jie ①甲酸酯：\ce{HCOOCH2C6H5}（1 种）、\ce{HCOOC6H4CH3}（甲基在邻、间、对，3 种）；②乙酸苯酯 \ce{CH3COOC6H5}（1 种）；③苯甲酸甲酯 \ce{C6H5COOCH3}（1 种）。共 6 种。\\
4 组峰：\ce{HCOOCH2C6H5}（苯环上 3 种 H $+$ \ce{CH2} $+$ HCOO 中的 H，共 5 组，不符）；甲酸对甲基苯酯 \ce{HCOO-C6H4-CH3}（对位）有 HCOO、苯环 2 种、\ce{CH3}，共 4 组，符合；\ce{CH3COOC6H5} 有 \ce{CH3} 和苯环 3 种，共 4 组，符合；\ce{C6H5COOCH3} 同理 4 组，符合。所以有 3 种。\\
\xianjing 一取代苯环上有邻、间、对 3 种 H；对位二取代（两基团不同）苯环上只有 2 种 H。
\end{liti}

\begin{liti}[共线共面]
\ce{CH3-CH=CH-C#C-CH3} 分子中，最多有几个碳原子共面？一定共线的碳原子有几个？
\tcblower
\jie 以 C=C 为中心：C1、C2、C3、C4 一定共面（乙烯模型）；C3—C4$\equiv$C5—C6 中，C4、C5 为 sp 杂化，C3、C4、C5、C6 一定共线，这条直线上的 C3、C4 已在平面内，所以 C5、C6 也在该平面内。6 个碳原子都共面；一定共线的是 C3、C4、C5、C6 共 4 个（C2—C3 不在这条直线上，因为 C3 是 $\mathrm{sp}^2$ 杂化，键角 120°）。
\end{liti}

\subsection{有机物的命名}
\textbf{烷烃系统命名}：选最长碳链为主链（等长时选支链多的），从离支链最近的一端编号（位次和最小），写成“位次—取代基名称—主链名称”，相同取代基合并并用“二、三”表示，不同取代基简单的写在前面。如 \ce{(CH3)2CHCH2C(CH3)3} 为 2,2,4-三甲基戊烷。\textbf{含官能团的有机物}：选含官能团的最长碳链，从离官能团最近的一端编号，标出官能团位置，如 \ce{CH2=CHCH(CH3)CH3} 为 3-甲基-1-丁烯，\ce{CH3CH(OH)CH3} 为 2-丙醇。苯的同系物可用邻、间、对或 1,2-二甲苯等表示。

\section{研究有机物的一般方法}
\textbf{分离提纯}：蒸馏（沸点相差较大的液体混合物）；重结晶（溶解度随温度变化差异大的固体，如粗苯甲酸：加热溶解$\to$趁热过滤$\to$冷却结晶$\to$过滤洗涤，2024年广东卷第6题、2025年第17题(2)均考查）；萃取分液；色谱法。

\textbf{确定实验式和分子式}：元素分析（燃烧法测 C、H，差值得 O）求实验式；\textbf{质谱}中最大质荷比即相对分子质量。

\textbf{确定结构}：\textbf{红外光谱}判断化学键和官能团；\textbf{核磁共振氢谱}中，吸收峰的\textbf{组数}$=$等效氢的种类数，峰面积之比$=$各类氢原子数之比。如乙醇 \ce{CH3CH2OH} 有 3 组峰，面积比 3:2:1；二甲醚 \ce{CH3OCH3} 只有 1 组峰。2025年广东卷第20题(4)：丙酮只有 1 组峰，溴代后生成 \ce{CH3COCH2Br}，单峰消失，出现两组新峰，面积比 3:2。

\section{烃的结构与性质}

\subsection{烷烃}
结构：碳原子均为 $\mathrm{sp}^3$ 杂化，链状烷烃通式 \ce{C_nH_{2n+2}}。性质稳定，不能使酸性 \ce{KMnO4} 溶液、溴水褪色；特征反应是光照下与卤素单质的\textbf{取代反应}（连锁反应，产物为多种卤代物的混合物）：\ce{CH4 + Cl2 ->[\text{光}] CH3Cl + HCl}。

\subsection{烯烃和炔烃}
\textbf{(1) 加成反应}：与 \ce{Br2}（溴水或溴的 \ce{CCl4} 溶液褪色）、\ce{H2}、HX、\ce{H2O} 等加成。不对称烯烃与 HX 加成时，主要产物中 H 加到含 H 较多的双键碳上（马氏规则），如 \ce{CH3CH=CH2 + HBr -> CH3CHBrCH3}（主要产物）。共轭二烯烃有 1,2-加成和 1,4-加成：\ce{CH2=CH-CH=CH2 + Br2 -> CH2Br-CH=CH-CH2Br}（1,4-加成）。\\
\textbf{(2) 氧化反应}：使酸性 \ce{KMnO4} 溶液褪色。\ce{RCH=CH2} 被氧化生成 \ce{RCOOH} 和 \ce{CO2}；\ce{R2C=CHR$'$} 生成酮 \ce{R2C=O} 和 \ce{R$'$COOH}——这一规律可用于由氧化产物反推双键位置。\\
\textbf{(3) 加聚反应}：\ce{$n$CH2=CH2 ->[\text{催化剂}] -[CH2-CH2]_{$n$}-}；\ce{$n$CH2=CHCl -> -[CH2-CHCl]_{$n$}-}（聚氯乙烯，单体是氯乙烯 \ce{CH2=CHCl}，不是氯乙烷——2025年广东卷第5题A）。\\
\textbf{(4) 乙炔的实验室制法}：\ce{CaC2 + 2H2O -> Ca(OH)2 + CH#CH ^}；用饱和食盐水代替水以减缓反应速率；用 \ce{CuSO4} 溶液除去 \ce{H2S} 等杂质。

\subsection{苯和苯的同系物}
\textbf{苯}：6 个碳原子均为 $\mathrm{sp}^2$ 杂化，形成大 $\pi$ 键，6 个碳碳键完全相同。苯不能使溴水（化学反应）和酸性 \ce{KMnO4} 溶液褪色；易取代、难加成、难氧化：
\begin{itemize}
\item 溴代：\ce{C6H6 + Br2 ->[\ce{FeBr3}] C6H5Br + HBr}（用液溴，不用溴水；产物溴苯因溶有 \ce{Br2} 呈褐色，用 NaOH 溶液洗涤后分液）。
\item 硝化：\ce{C6H6 + HNO3 ->[\text{浓硫酸}][\text{50～60 ℃}] C6H5NO2 + H2O}（水浴加热；浓硫酸作催化剂和吸水剂）。
\item 加成：\ce{C6H6 + 3H2 ->[\ce{Ni}][\text{△}] C6H12}（环己烷）。
\end{itemize}

\textbf{苯的同系物}（如甲苯）——\textbf{基团间相互影响}：①苯环使侧链活化：与苯环直接相连的碳上有 H 时，侧链能被酸性 \ce{KMnO4} 溶液氧化为 \ce{-COOH}（如甲苯$\to$苯甲酸，2025年广东卷第17题(1)现象为“溶液由紫色变为无色”）；②侧链使苯环活化：甲苯硝化在邻、对位发生，可生成 2,4,6-三硝基甲苯（TNT）。

\begin{yicuo}
\begin{enumerate}
\item 苯与溴水混合振荡，溴水褪色是\textbf{萃取}（物理变化），不是加成。
\item 苯的同系物侧链被 \ce{KMnO4} 氧化的前提是“与苯环直接相连的碳原子上有氢原子”，如叔丁基苯不能被氧化。
\item 光照下甲苯与 \ce{Cl2} 取代侧链上的 H；\ce{FeCl3} 催化下取代苯环上的 H——条件不同，位置不同。
\end{enumerate}
\end{yicuo}

\begin{center}
\small
\begin{tabularx}{\linewidth}{p{1.8cm}p{1.9cm}p{2.0cm}X}
\toprule
烃 & 溴水（或 \ce{Br2}/\ce{CCl4}） & 酸性 \ce{KMnO4} & 特征反应 \\
\midrule
烷烃 & 不褪色 & 不褪色 & 光照取代 \\
烯烃、炔烃 & 褪色（加成） & 褪色（氧化） & 加成、氧化、加聚 \\
苯 & 不反应（萃取溴水中的溴，水层褪色；与 \ce{Br2}/\ce{CCl4} 互溶，不褪色） & 不褪色 & \ce{FeBr3} 催化溴代、硝化、加氢 \\
甲苯等 & 不反应（同苯，只发生萃取） & 褪色（侧链氧化） & 侧链氧化；苯环邻、对位取代；光照侧链取代 \\
\bottomrule
\end{tabularx}
\end{center}

\begin{fangfa}[烃的燃烧规律]
\ce{C_xH_y + ($x$ + $\frac y4$)O2 -> $x$CO2 + $\frac y2$H2O}。①等物质的量的烃完全燃烧，耗氧量看 $x+\frac y4$；②等质量的烃完全燃烧，H 的质量分数越大，耗氧越多（\ce{CH4} 最多）；③120 ℃ 以上（水为气体）燃烧前后气体体积不变时，$y=4$（如 \ce{CH4}、\ce{C2H4}、\ce{C3H4}）。
\end{fangfa}

% ---------------- 章末 ----------------
\zhangmo{本章小结}
\subsection*{一、知识总结}
\begin{center}
\small
\begin{tabularx}{\linewidth}{p{2.2cm}X}
\toprule
结构 & 饱和碳 $\mathrm{sp}^3$、双键碳 $\mathrm{sp}^2$、三键碳 sp；甲烷、乙烯、乙炔、苯四种空间模型 \\
同分异构 & 碳架、位置、官能团、顺反、对映；等效氢、基元法、定一移一、苯环取代 3/6/10 \\
研究方法 & 蒸馏、重结晶、萃取、色谱；质谱（相对分子质量）、红外（官能团）、核磁（氢的种类与数目比） \\
烷烃 & 光照取代 \\
烯、炔 & 加成（\ce{Br2}、\ce{H2}、HX、\ce{H2O}）、氧化（\ce{KMnO4} 褪色）、加聚 \\
芳香烃 & 苯：溴代（\ce{FeBr3}）、硝化、加成；甲苯：侧链氧化、邻对位取代 \\
\bottomrule
\end{tabularx}
\end{center}
\subsection*{二、方法提升}
\begin{enumerate}
\item 同分异构体书写“三步走”：先算不饱和度、定类别（官能团），再定碳架，最后移位置；用对称性去重。
\item 用核磁共振氢谱“峰组数$+$峰面积比”验证结构；限定条件题先把“必须有的结构”拿出来，剩余部分再组合。
\end{enumerate}
\subsection*{三、易错点归纳}
\begin{enumerate}
\item 认为含 \ce{-CH3} 的分子中所有原子可能共面（错）。
\item 苯环上二元取代计数遗漏或重复；“两同一异”记成 10 种。
\item 写加聚产物时把双键碳以外的原子放进主链。
\item 苯与溴水“褪色”误认为加成反应。
\end{enumerate}

\zhangmo{分层练习}
\zu{A组\quad 基础巩固}
\begin{xiti}[start=1]
\item 下列有机物分子中所有原子一定共平面的是\xz
\xxb{\ce{CH3CH=CH2}}{苯乙烯}{\ce{CH2=CH-CH=CH2}}{\ce{C6H5Cl}（氯苯）}

\item 分子式为 \ce{C4H8O2} 且属于酯的同分异构体有\tk{}种；分子式为 \ce{C5H10O2} 且能与 \ce{NaHCO3} 反应放出气体的有\tk{}种。

\item 某烃的相对分子质量为 72，其核磁共振氢谱只有 1 组峰，则其结构简式为\tk，系统命名为\tk。

\item 下列关于苯及其同系物的说法正确的是\xz
\xxa{苯能使溴水褪色，发生了加成反应}{甲苯能使酸性 \ce{KMnO4} 溶液褪色，说明苯环对甲基有影响}{苯与液溴在光照条件下生成溴苯}{苯的硝化反应中，浓硫酸只作吸水剂}

\item 用系统命名法命名：(1) \ce{CH2=C(CH3)-CH=CH2}\tk；(2) \ce{CH3CH(C2H5)CH2CH3}\tk；(3) \ce{CH#C-CH(CH3)2}\tk。

\item 等质量的 \ce{CH4}、\ce{C2H4}、\ce{C2H2}、\ce{C6H6} 完全燃烧，消耗 \ce{O2} 最多的是\tk；等物质的量的上述四种烃完全燃烧，消耗 \ce{O2} 最多的是\tk。
\end{xiti}

\zu{B组\quad 高考典型}
\begin{xiti}
\item \laiyuan{2024年广东卷第20题(2)，来源：仓库所附真题文档}\ 化合物Ⅱ为苯乙酮（\ce{C6H5COCH3}），其中含氧官能团的名称是\tk。化合物Ⅱ的某同分异构体含有苯环，在核磁共振氢谱上只有 4 组峰，且能够发生银镜反应，其结构简式为\tk。

\item \laiyuan{原创}\ 某烯烃 A（\ce{C6H12}）被酸性 \ce{KMnO4} 溶液氧化，只得到一种产物——丙酮（\ce{CH3COCH3}）。(1) 写出 A 的结构简式并命名。(2) A 与 HBr 加成的产物有几种？(3) A 是否存在顺反异构？

\item \laiyuan{原创}\ 分子式为 \ce{C8H10} 的芳香烃共有\tk{}种；其中苯环上的一氯代物只有一种的是\tk（写结构简式）；被酸性 \ce{KMnO4} 溶液氧化后生成苯甲酸的是\tk（写名称）。

\item \laiyuan{原创}\ 分子式为 \ce{C5H10} 的链状烯烃，不考虑顺反异构时有\tk{}种；考虑顺反异构时有\tk{}种。其中与 \ce{H2} 加成后得到 2-甲基丁烷的有\tk{}种。

\item \laiyuan{原创}\ 分子式为 \ce{C9H12} 的芳香烃共有\tk{}种；其中苯环上的一溴代物只有一种的结构简式为\tk。
\end{xiti}

\zu{C组\quad 冲刺高分}
\begin{xiti}
\item \laiyuan{原创}\ 某有机物 X 只含 C、H、O 三种元素。取 1.48 g X 完全燃烧，生成 3.52 g \ce{CO2} 和 1.80 g \ce{H2O}；质谱显示最大质荷比为 74；红外光谱显示有 O—H 键；核磁共振氢谱有 2 组峰，面积比为 9:1。(1) 求 X 的分子式。(2) 写出 X 的结构简式并命名。(3) 写出与 X 同类别（醇）的同分异构体中，能被催化氧化为醛的结构简式。

\item \laiyuan{原创}\ 0.1 mol 烃 A 完全燃烧生成 0.6 mol \ce{CO2} 和 0.5 mol \ce{H2O}；A 能使溴的 \ce{CCl4} 溶液褪色，1 mol A 最多与 1 mol \ce{H2} 加成；A 被酸性 \ce{KMnO4} 溶液氧化只得到己二酸 \ce{HOOC(CH2)4COOH}。(1) 求 A 的分子式并计算不饱和度。(2) 写出 A 的结构简式（名称）。(3) 写出 A 与 \ce{Br2} 反应的化学方程式。(4) A 的同分异构体中，属于链状二烯烃且核磁共振氢谱只有 2 组峰的结构简式为\tk。
\end{xiti}

\zhangmo{练习解析}
\begin{jiexilist}
\item \ans{D}A 含 \ce{-CH3}，不可能全部共面；B、C 中苯环与乙烯基、两个乙烯基之间以单键相连，可旋转，“可能”共面但不一定；D 中 Cl 取代苯环上的 H，仍在苯环平面内，一定共面。
\item 4；4。酯：甲酸正丙酯、甲酸异丙酯、乙酸乙酯、丙酸甲酯。羧酸 \ce{C4H9COOH}：丁基 4 种。
\item \ce{C(CH3)4}；2,2-二甲基丙烷。$72=14\times5+2$，符合 \ce{C_nH_{2n+2}}（$n=5$），为 \ce{C5H12}；只有一种等效氢，为新戊烷。
\item \ans{B}A：萃取；B：甲基本身不能被 \ce{KMnO4} 氧化，在苯环影响下才被氧化为羧基，正确；C：需 \ce{FeBr3} 催化，不是光照；D：浓硫酸作催化剂和吸水剂。
\item (1) 2-甲基-1,3-丁二烯（异戊二烯）；(2) 3-甲基戊烷（主链应选含乙基的 5 个碳，不能命名为“2-乙基丁烷”）；(3) 3-甲基-1-丁炔。
\item \ce{CH4}；\ce{C6H6}。等质量时含 H 质量分数越大耗氧越多（\ce{CH4} 最大）；等物质的量时看 $x+\frac y4$：\ce{CH4} 2、\ce{C2H4} 3、\ce{C2H2} 2.5、\ce{C6H6} 7.5。
\item 酮羰基；对甲基苯甲醛（\ce{CH3C6H4CHO}，对位）。能发生银镜反应说明含 \ce{-CHO}；\ce{C8H8O} 除苯环和醛基外剩 1 个 C，为甲基；4 组峰要求对称性高，甲基与醛基处于对位（氢的种类：\ce{-CH3}、苯环上两种、\ce{-CHO}）。
\item (1) \ce{(CH3)2C=C(CH3)2}，2,3-二甲基-2-丁烯。氧化只得丙酮，说明双键两端各连两个甲基。(2) 1 种。分子结构对称，与 HBr 加成只得一种产物 \ce{(CH3)2CBr-CH(CH3)2}。(3) 不存在，双键碳上各连两个相同的甲基。
\item 4；\ce{CH3-C6H4-CH3}（对二甲苯）；乙苯。\ce{C8H10} 芳香烃有乙苯和邻、间、对二甲苯 4 种。对二甲苯苯环上的 4 个 H 等效，一氯代物只有 1 种。侧链被酸性 \ce{KMnO4} 氧化时，无论侧链多长，与苯环相连的碳都变为 \ce{-COOH}：乙苯（1 个侧链）生成苯甲酸，二甲苯（2 个侧链）生成苯二甲酸。
\item 5；6；3。碳架 C—C—C—C—C：1-戊烯、2-戊烯（有顺反异构）；碳架 C—C(C)—C—C：2-甲基-1-丁烯、3-甲基-1-丁烯、2-甲基-2-丁烯；新戊烷碳架不能形成双键。加氢得 2-甲基丁烷的是带支链的 3 种。
\item 8；1,3,5-三甲苯（均三甲苯），结构简式如下：\par{\centering\setchemfig{atom sep=1.5em}\chemfig{*6(-(-CH_3)=-(-CH_3)=-(-CH_3)=)}\par}一个侧链：丙苯、异丙苯（2 种）；两个侧链（甲基、乙基）：邻、间、对（3 种）；三个甲基：连三甲苯、偏三甲苯、均三甲苯（3 种）。均三甲苯高度对称，苯环上 3 个 H 等效。
\item (1) $n(\ce{C})=3.52/44=0.08$ mol，$n(\ce{H})=2\times1.80/18=0.20$ mol，$m(\ce{O})=1.48-0.96-0.20=0.32$ g，$n(\ce{O})=0.02$ mol；C:H:O$=4:10:1$，实验式 \ce{C4H10O}，式量 74，分子式 \ce{C4H10O}。(2) 有 O—H 键，为醇；2 组峰面积比 9:1，为 \ce{(CH3)3COH}，2-甲基-2-丙醇。(3) 能氧化为醛的醇含 \ce{-CH2OH}：\ce{CH3CH2CH2CH2OH}、\ce{(CH3)2CHCH2OH}。
\item (1) \ce{C6H10}，$\Omega=\frac{2\times6+2-10}{2}=2$。(2) 只与 1 mol \ce{H2} 加成，说明只有 1 个碳碳双键，另 1 个不饱和度为环；氧化只得到直链二元酸，双键在六元环中：环己烯。(3) 加成反应，生成 1,2-二溴环己烷：\par{\centering\setchemfig{atom sep=1.5em}\chemfig{*6(-=----)}\quad\ce{+ Br2 ->}\quad\chemfig{*6(-(-Br)-(-Br)----)}\par}(4) \ce{CH2=C(CH3)-C(CH3)=CH2}（2,3-二甲基-1,3-丁二烯，\ce{CH2} 与 \ce{CH3} 两种 H，面积比 4:6）。
\end{jiexilist}
